\section{Calibration Model Validation}
\label{app:calibration-validation}
\label{app:forensics}

This appendix defines the diagnostics that make a calibrated CJE estimate
credible. It does not duplicate the result tables. Each diagnostic answers a
specific design question.

\subsection{Signal in the Cheap Score}
\label{app:reliability}

The cheap score must contain usable information about the oracle outcome,
especially in the lower tail. The basic diagnostic is a decile table:
\[
\text{bin }b
\quad\mapsto\quad
\left(
\overline{\Sscore}_b,\;
\overline{\Y}_b,\;
\overline{\Sscore}_b-\overline{\Y}_b,\;
\mathrm{sd}(\Y\mid b)
\right).
\]
The table is not a proof of transport. It checks that the surrogate has signal
and that the cheap-to-oracle map is not flat in the region where CVaR is
estimated.

\subsection{Tail-Specific Calibration Gap}

Mean calibration can look good while the tail is wrong. For each policy and
tail level, report
\[
\CVaR_\alpha(\Sscore),\qquad
\CVaR_\alpha(\Y),\qquad
\CVaR_\alpha(\Y)-\CVaR_\alpha(\Sscore).
\]
The purpose is to show whether the cheap judge compresses, misses, or reverses
tail severity before calibration. This diagnostic motivates the stop-loss
calibrator but does not gate the level claim; the audit does.

\subsection{Base-vs-Clone Sanity Check}
\label{app:bcforensics}

When two policies share model, prompt, and sampling configuration except for
seed, their distributions should be statistically close. A large calibrated
gap between them is evidence of pipeline or grader instability. The diagnostic
should report:
\begin{itemize}[leftmargin=*,nosep]
\item mean difference;
\item CVaR difference;
\item overlap of lower-tail prompt identities;
\item whether disagreements concentrate in a small set of rubric criteria.
\end{itemize}

\subsection{Policy-Wise Residual Transport}

For each estimand, validation is policy-wise:
\[
\text{mean:}\quad \Y-\hat f(\Sscore,\X),
\qquad
\text{CVaR:}\quad
\hat g_1,\hat g_2.
\]
The diagnostic should show residual means and audit verdicts by policy. It
should also state the refusal rule: a rejected row is a diagnostic point
estimate, not an interpretable calibrated level.

\subsection{Risky-Policy Mechanism Check}

For a policy designed to expose mean-vs-tail divergence, the key validation is
not only the final mean and CVaR. It is whether the failure concentrates where
the policy is designed to fail. Split prompts into routine and dangerous
medical edge cases and report:
\[
\overline{\Y}_{\mathrm{routine}},
\quad
\overline{\Y}_{\mathrm{dangerous}},
\quad
\CVaR_{\alpha,\mathrm{routine}},
\quad
\CVaR_{\alpha,\mathrm{dangerous}}.
\]
The desired pattern is ordinary routine performance and a much worse dangerous
tail. Without this split, a constructed risky policy can look like a tuned row
rather than a mechanism-backed stress test.

\subsection{Threshold-Gap Diagnostic}
\label{app:thresholdgap}

Let $\hat t_\alpha$ be the threshold selected by the calibrated surrogate and
let $t^\star_\alpha$ be the atom-split oracle quantile on the full oracle panel
when that panel is available. The gap
\[
\hat t_\alpha-t^\star_\alpha
\]
is a mechanism diagnostic. It does not gate the level claim. A large gap with
small audit moments can happen because the stop-loss objective trades off tail
mass and shortfall severity. A large gap with large audit moments is evidence
that the threshold mechanism itself failed to transport.

\subsection{Implementation Invariants From the Math Notes}

The estimator should satisfy these non-negotiable checks:
\begin{itemize}[leftmargin=*,nosep]
\item the threshold grid covers the upper endpoint needed for the
$\alpha=1$ identity;
\item the stop-loss calibrator is monotone decreasing in the cheap score;
\item the mean calibrator is monotone increasing in the cheap score;
\item bootstrap replicates refit the calibrator and re-optimize
$\hat t_\alpha$;
\item at $\alpha=1$, CVaR truth, point estimates, bootstrap draws,
calibration variance, and audit residuals collapse to their Mean-CJE
counterparts.
\end{itemize}
These are regression-test targets, not optional diagnostics.
